% Lecture visuelle directe avec la vision du modèle ; aucun OCR.
% Source : PDF original, vue 038, côté droit, page manuscrite 77.
% Lot p. 77–83 ; chapitre à poursuivre p. 84.
% Statut : transcrit ; relecture indépendante à effectuer.
\part*{Histoire de la Congrégation des jeunes gens depuis sa fondation jusqu'à celle de la Congrégation des hommes\\[1em]{\normalfont\large\textit{Ad majorem Dei gloriam}}\\[0.5em]{\normalfont\small Pour la plus grande gloire de Dieu. (Note du transcripteur)}}
\addcontentsline{toc}{part}{Histoire de la Congrégation depuis sa fondation}
\setcounter{chapter}{0}
\renewcommand{\theHchapter}{histoire.\arabic{chapter}}
\chapter{Fondation de la Congrégation}

\section*{Résolution de fonder la Congrégation}
% Mention marginale de la page manuscrite 77.
Dieu avoit tout préparé d'avance pour l'exécution des desseins de miséricorde qu'il avoit conçus sur nous. L'époque est enfin arrivée où nous allons en voir l'accomplissement. Les circonstances ne pouvoient pas être plus favorables ; le rétablissement de la religion avançoit à grands pas. À la vue de tout ce qui se passoit, M. Delagrange étoit trop zélé pour rester indécis et nous laisser temporiser encore. Il va trouver d'une part M. Ripoud, de l'autre l'auteur de cette histoire et les somme de tenir enfin la parole qu'ils avoient donnée de mettre la main à l'œuvre pour l'établissement de la Congrégation aussitôt que le concordat seroit à Lyon en voie d'exécution. Il n'y auroit eu alors aucun motif raisonnable de retarder davantage ; nous consentîmes donc tous les deux à entrer en conférence
% Source : PDF original, vue 039, côté gauche, page manuscrite 78.
pour donner suite à nos projets. Dans ce premier colloque entre nous trois, deux points essentiels furent agités ; nous comprîmes la nécessité qu'il y avoit pour nous de choisir un prêtre animé de l'esprit de Dieu qui voulût bien se charger de la direction de la Congrégation ; avant d'aller en avant nous crûmes donc convenable de déterminer ce choix. C'étoit ce qui nous paroissoit le plus pressé. Plusieurs avis furent ouverts ; pour mon compte, en rappelant les premières tentatives qui avoient été faites dans le temps pour le même objet sous les auspices de M. de Villers, je connoissois tant le zèle de ce digne ecclésiastique et son aptitude pour conduire les cœurs et les âmes. J'aurois donc désiré lui confier le soin de la Congrégation naissante. Mais ces Messieurs, tout en rendant justice au mérite de M. de Villers qu'ils connoissoient aussi bien que moi, m'objectèrent qu'il y auroit encore plus d'avantages pour nous si nous pouvions réussir à nous mettre sous le patronage d'un ordre religieux. Nous avons dans notre ville les Pères de la foi ; en nous unissant à eux, en nous greffant pour ainsi dire sur cette corporation qui avoit déjà acquis une bien grande considération, nous trouverons là un appui constant, et en cas de changement de directeur nous serons assurés que la direction demeurera toujours la même. Je n'avois rien à répondre à des raisons aussi solides ; je m'y rendis donc et tout d'une voix nous convînmes de prier les Pères de la foi de se charger de nous. Le moment étoit opportun ; le P. Varin, leur provincial, se trouvoit pour lors à Lyon. Nous
% Source : PDF original, vue 039, côté droit, page manuscrite 79.
n'hésitons plus, nous décidons d'aller lui faire notre confidence et de lui présenter notre requête. Avant de faire cette démarche il nous restoit encore un autre point à traiter. Jusqu'ici nous n'étions que trois ; il nous sembla que même pour commencer ce n'étoit pas assez. Nous songeâmes donc à nous adjoindre de nouveaux frères. M. Delagrange proposa M. Lallemand ; je proposai MM. Bori et Journoud, tous trois furent agréés et sur les propositions qui leur en fut faite ils acceptèrent avec empressement\footnote{Raccord de « sur les propositions qui leur en fut faite » à confirmer ; noms lus sur le scan, Q040. (Note du transcripteur).}. Mais avant même de les avoir mis dans notre confidence, aussitôt après la résolution de notre petit colloque, nous nous hâtons de nous rendre dans la maison où les Pères de la foi vivent en communauté pour savoir sur quoi nous pouvions compter à l'égard de la demande que nous avions à leur faire. Nous y rencontrons le P. Varin ; à côté de lui se trouvoit le P. Roger. Nous exposons au P. Varin nos pensées et nos vœux ; il nous accueille avec bonté, nous écoute avec intérêt et se retournant du côté du P. Roger il nous le montre en nous disant : Messieurs, voilà votre homme. À ces paroles le visage du P. Roger s'illumine ; à voir l'expression de sa physionomie il étoit impossible de douter que son âme toute de charité avoit compris l'importance de la mission qui lui étoit confiée et s'en chargeoit avec joie. Nous vîmes à l'instant dans ses yeux qu'il étoit devenu notre père et que nous pouvions déjà nous considérer comme ses enfants. Nous prîmes jour au moment même avec lui pour fixer dans une conférence les
% Source : PDF original, vue 040, côté gauche, page manuscrite 80.
bases principales sur lesquelles devoit s'élever l'édifice de la Congrégation et arrêter l'époque définitive de sa fondation. Nous choisîmes pour cette réunion le 14 juillet, le jour où devoit se célébrer encore cette année (et on célébra l'année suivante pour la dernière fois) la fête de la fédération. Quoique ce ne fût certainement pas une fête pour nous, elle nous empêchoit néanmoins de nous livrer à nos occupations ordinaires et nous donnoit par force un peu de loisir que nous ne pouvions pas mieux employer. Singulier rapprochement que la Providence permettoit ici. Cette fédération de la révolte alloit succomber sous le joug du despotisme, héritier de la Révolution ; mais auparavant il faut qu'elle cède la place à une nouvelle fédération toute chrétienne ; celle-là n'enfantera pas la liberté ni la licence\footnote{Lecture du premier terme de l'opposition à confirmer, Q041. (Note du transcripteur).}, mais elle assurera à ceux qui y seront admis cette véritable liberté qui est le partage des enfants de Dieu, et elle les réunira dans le but de travailler à rendre au Seigneur les âmes que la première lui a enlevées.

\section*{Conférence du 14 juillet}
% Mention marginale de la page manuscrite 80.
Les trois amis que nous avions choisis pour fonder avec nous la Congrégation furent avertis de se trouver à ce rendez-vous. Tous les six s'y rendirent exactement. La conférence dura environ deux heures. Les articles préliminaires qui devoient servir de base au règlement furent arrêtés d'un commun accord avec le P. Roger. Il fut décidé dans cette assemblée : 1\up{o} que le but de la Congrégation seroit la gloire de Dieu, l'honneur de Marie, notre propre sanctification et celle du prochain ; 2\up{o} qu'elle auroit la Ste Vierge pour patronne ; 3\up{o} que la
% Source : PDF original, vue 040, côté droit, page manuscrite 81.
fondation auroit lieu le dimanche suivant 18 juillet ; 4\up{o} qu'on s'occuperoit après la fondation de rédiger un règlement sur les bases qui venoient d'être arrêtées ; 5\up{o} on admit M. Jordan comme fondateur, mais on décida en même temps que leur nombre étoit irrévocablement fixé à sept et qu'en conséquence on ne proposeroit plus personne.

\section*{Fondation}
% Mention marginale de la page manuscrite 81.
Ce jour mémorable qui a laissé parmi nous de si profonds souvenirs, le 3\up{e} dimanche de juillet 1802, le 18\up{e} jour du mois paroît enfin. Dès cinq heures du matin les sept fondateurs dont je vais vous répéter les noms se réunissent dans le petit appartement bien simple, bien modeste qu'habitoient dans le quartier St Paul les Pères de la foi. Ces fondateurs étoient MM. Ripoud, Delagrange, Lallemand, Journoud, Bori, Jordan et Coste. Ils se placent à côté les uns des autres au pied du modeste autel élevé dans la chapelle des Pères. Le P. Roger entonne le \textit{Veni Creator}\footnote{Début de l'hymne latine : « Venez, Créateur. » (Note du transcripteur).} ; les sept amis, tous un cierge à la main, unissent leurs prières pour supplier l'Esprit-Saint, cet esprit d'amour et de sagesse, d'embraser leurs cœurs du feu de la charité afin que ce feu sacré pût se perpétuer constamment dans la société qui alloit se former et qu'il embrasât par la suite tous ceux qui viendroient se joindre à eux ; ils le supplient aussi de répandre dans leurs cœurs cette sagesse ineffable, vraie colonne lumineuse qui seule pouvoit éclairer la marche de la tribu chérie que Dieu daignoit assembler sous les auspices de Marie ; à l'offertoire le pain fut offert en commun à la bénédiction par les sept nouveaux frères ; il devenoit un symbole de cet autre pain qui en unissant
% Source : PDF original, vue 041, côté gauche, page manuscrite 82.
leurs sept cœurs dans le cœur du divin maître devoit cimenter à jamais l'alliance qu'ils alloient contracter. Avant la communion ils prononcent successivement en présence de l'adorable Eucharistie cette consécration qu'ils faisoient avec tant de bonheur de tout leur être à la mère de Dieu. Enfin J.-C. vient habiter en eux ; il vient apposer le sceau de son amour à toutes ces bonnes résolutions, tous ces pieux projets, tous ces saints désirs qu'il avoit inspirés lui-même. Lisez, Messieurs, Dieu seul connoît ce qui se passoit alors dans nos cœurs. Comme ils étoient pénétrés de l'ardent désir d'être tout à lui, d'employer tous ces dons que nous avions reçus de sa bonté infinie à procurer sa gloire ! Comme la charité les remplissoit ! Comme elle nous faisoit chérir par avance d'un amour de dilection tous ces frères que nous apercevions par la pensée accourir d'années en années pour se joindre à nous ! Comme nous les présentions à Marie ! Comme nous la suppliions de nous adopter, de les adopter tous avec nous ! Oui, je crois pouvoir vous le dire : tous les congréganistes qui depuis cette époque heureuse sont venus s'unir à nous, tous ceux qui viendront encore, tous ceux qui nous suivront tant qu'il plaira au Seigneur de laisser subsister cette sainte société, tous sans exception étoient en ce moment solennel l'objet de nos prières et de notre amour. Un \textit{Te Deum}\footnote{Début de l'hymne latine : « Vous, Dieu, [nous vous louons] ». (Note du transcripteur).} chanté avec enthousiasme ne fut encore qu'une foible expression du bonheur dont nos âmes surabondoient et des sentimens
% Source : PDF original, vue 041, côté droit, page manuscrite 83.
de reconnoissance dont elles étoient agitées.

Nous passâmes ensuite dans le réfectoire des Pères. Notre premier acte fut de mettre en réserve pour être conservé pour la Congrégation comme mémorial de cette douce union une portion de notre pain bénit ; des morceaux en furent également distribués à tous les fondateurs pour être aussi conservés par eux dans le même but. Vous savez qu'il n'y a pas encore bien longtemps à chaque anniversaire de la fondation ce pain vénérable empreint de si précieux souvenirs reparoissoit au milieu de nous et qu'on en distribuoit des parcelles à chaque congréganiste pour les rattacher en quelque manière par la participation à la même nourriture à tous les temps divers de la Congrégation. Le déjeuner fut égayé par les aimables saillies d'une joie bien pure. Le P. Roger avoit un admirable talent pour nous en faire savourer les douceurs. Au milieu de ces délicieuses émotions dont nous étions pénétrés et de toutes ces effusions de la tendresse fraternelle, voilà que tout à coup nos bras s'entrelaçant le bon père est placé au centre ; tous les corps se trouvent réunis autour de lui comme l'étoient déjà tous les cœurs. Le père et les sept frères groupés ainsi ne font vraiment plus qu'un, c'est un faisceau d'amour que vient de former la charité. Là sous les ailes protectrices de notre nouveau père nous échangeons tous ensemble les plus doux témoignages de la tendre amitié qui nous unit à jamais : \textit{Sicut novellæ olivarum in circuitu mensæ tuæ}\footnote{Traduction littérale proposée : « Comme de jeunes oliviers autour de votre table. » Graphie latine et référence à vérifier, Q045. (Note du transcripteur).}. C'est ainsi, Messieurs, que se donna pour la première fois dans la Congrégation le baiser de charité.
